% Problem-section file following ../_template.tex.
\section{Exactly solvable nondegradable quantum channels}
\label{sec:problem-74}

\paragraph{Problem.}
Does there exist a finite-dimensional quantum channel that is neither
degradable nor antidegradable and whose unassisted quantum capacity is known
exactly?  For a channel $\mathcal N_{A\to B}$ with complementary channel
$\mathcal N^c_{A\to E}$, degradability and antidegradability mean,
respectively, that there is a channel $\mathcal D$ or $\mathcal A$ satisfying
\begin{equation}
  \mathcal N^c=\mathcal D_{B\to E}\circ\mathcal N,
  \qquad\text{or}\qquad
  \mathcal N=\mathcal A_{E\to B}\circ\mathcal N^c.
  \label{eq:p74-degradability-definitions}
\end{equation}
The channel sought must satisfy neither identity in
Eq.~\eqref{eq:p74-degradability-definitions}.  Its quantum capacity is
defined by the coherent-information regularization
\begin{equation}
  Q(\mathcal N)
  :=\sup_{n\geq1}\frac1n\max_{\rho_{A^n}}
  \left[
    S(\mathcal N^{\otimes n}(\rho_{A^n}))
    -S((\mathcal N^c)^{\otimes n}(\rho_{A^n}))
  \right].
  \label{eq:p74-quantum-capacity}
\end{equation}
The question asks for an explicit channel outside both classes in
Eq.~\eqref{eq:p74-degradability-definitions} together with an exact evaluation
of Eq.~\eqref{eq:p74-quantum-capacity}.

\subsection*{Status}

Solved

\subsection*{Source}

Wilde explicitly identified finding the quantum capacity of a nondegradable
channel as an important open challenge
\sourcecite{ref:p74-wilde}{Wil17}.

\subsection*{Progress}

\begin{itemize}
  \item Chessa and Giovannetti supplied a simple explicit qutrit solution to
  the existential problem.  For $1/2<\gamma<1$, define a channel
  $\mathcal N_\gamma$ by the Stinespring isometry
  \begin{equation}
    \begin{aligned}
      V_\gamma|0\rangle_A&=|0\rangle_B|0\rangle_E,\\
      V_\gamma|1\rangle_A
        &=\sqrt{1-\gamma}\,|1\rangle_B|0\rangle_E
          +\sqrt{\gamma}\,|0\rangle_B|1\rangle_E,\\
      V_\gamma|2\rangle_A&=|2\rangle_B|0\rangle_E,
      \qquad
      \mathcal N_\gamma(\rho):=
        \operatorname{Tr}_E[V_\gamma\rho V_\gamma^\dagger].
    \end{aligned}
    \label{eq:p74-single-decay-mad-channel}
  \end{equation}
  They proved that the channel in
  Eq.~\eqref{eq:p74-single-decay-mad-channel} is neither degradable nor
  antidegradable throughout this parameter range and nevertheless has
  \begin{equation}
    Q(\mathcal N_\gamma)=1.
    \label{eq:p74-mad-capacity}
  \end{equation}
  The noiseless subspace $\operatorname{span}\{|0\rangle,|2\rangle\}$ gives
  the achievable qubit in Eq.~\eqref{eq:p74-mad-capacity}; their channel
  analysis supplies the matching upper bound
  \sourcecite{ref:p74-chessa-giovannetti}{CG21}.

  \item Earlier work of D'Arrigo, Benenti, Falci, and Macchiavello determined
  the quantum capacity and studied the degradability properties of a fully
  correlated two-qubit amplitude-damping channel.  It is a historical
  antecedent to the multilevel constructions, but the qutrit witness in
  Eq.~\eqref{eq:p74-single-decay-mad-channel} independently suffices for the
  existential question here
  \sourcecite{ref:p74-darrigo-et-al}{DBF+13}.

  \item Smith and Wu later constructed broader families of genuinely
  nondegradable channels with additive coherent information, hence with
  quantum capacities given by a one-use optimization.  Their results provide
  further exactly solvable examples but postdate the family in
  Eq.~\eqref{eq:p74-single-decay-mad-channel}
  \sourcecite{ref:p74-smith-wu}{SW25}.
\end{itemize}

\subsection*{References}

\begin{enumerate}[label={},leftmargin=4.8em,labelsep=0.5em]
  \footnotesize
  \item[\textup{[Wil17]}]\label{ref:p74-wilde}
  M. M. Wilde, \emph{Quantum Information Theory}, 2nd ed., Cambridge
  University Press (2017), Sec.~26.6.
  \href{https://doi.org/10.1017/9781316809976}%
       {doi:10.1017/9781316809976};
  \href{https://arxiv.org/abs/1106.1445}{arXiv:1106.1445}.

  \item[\textup{[CG21]}]\label{ref:p74-chessa-giovannetti}
  S. Chessa and V. Giovannetti, ``Quantum Capacity Analysis of Multi-Level
  Amplitude Damping Channels,'' \emph{Communications Physics} \textbf{4}, 22
  (2021).
  \href{https://doi.org/10.1038/s42005-021-00524-4}%
       {doi:10.1038/s42005-021-00524-4};
  \href{https://arxiv.org/abs/2008.00477v3}{arXiv:2008.00477v3}.

  \item[\textup{[DBF+13]}]\label{ref:p74-darrigo-et-al}
  A. D'Arrigo, G. Benenti, G. Falci, and C. Macchiavello,
  ``Classical and Quantum Capacities of a Fully Correlated Amplitude Damping
  Channel,'' \emph{Physical Review A} \textbf{88}, 042337 (2013).
  \href{https://doi.org/10.1103/PhysRevA.88.042337}%
       {doi:10.1103/PhysRevA.88.042337};
  \href{https://arxiv.org/abs/1309.2219}{arXiv:1309.2219}.

  \item[\textup{[SW25]}]\label{ref:p74-smith-wu}
  G. Smith and P. Wu, ``Additivity of Quantum Capacities in Simple
  Non-Degradable Quantum Channels,'' \emph{IEEE Transactions on Information
  Theory} \textbf{71}, 6134--6154 (2025).
  \href{https://doi.org/10.1109/TIT.2025.3583936}%
       {doi:10.1109/TIT.2025.3583936};
  \href{https://arxiv.org/abs/2409.03927v4}{arXiv:2409.03927v4}.
\end{enumerate}

\subsection*{Comment}

The family in Eq.~\eqref{eq:p74-single-decay-mad-channel} solves the literal
existence problem.  It does not provide a general single-letter formula or a
classification of nondegradable channels with additive coherent information.

\subsection*{Tag}

Amplitude-damping channels; Quantum capacity; Coherent information; Quantum
channels; Qutrit systems; Quantum Shannon theory

\subsection*{ID}

\texttt{op\_a59d7cc1c843edb5}
